\documentclass[12pt,a4paper]{article}
\usepackage[UTF8]{ctex}
\usepackage{graphicx}
\usepackage{geometry}
\usepackage{amsmath}
\usepackage{amsfonts}
\usepackage{amssymb}
\usepackage{listings}
\usepackage{xcolor}
\usepackage{float}
\usepackage{tikz}
\usepackage{fancyhdr}
\usepackage{setspace}
\usepackage{multirow}
\usepackage{booktabs}
\usepackage[most]{tcolorbox}
\tcbset{colback=white,colframe=black,boxrule=0.6pt,arc=0mm,left=3mm,right=3mm,top=2mm,bottom=2mm,breakable,enhanced}

\newcommand{\fillin}[1]{\underline{\makebox[8cm][l]{#1}}}
\usetikzlibrary{shapes,arrows,positioning,calc,chains}

% 确保section标题使用宋体
\makeatletter
\renewcommand\section{\@startsection{section}{1}{\z@}%
    {-3.5ex \@plus -1ex \@minus -.2ex}%
    {2.3ex \@plus.2ex}%
    {\normalfont\Large\bfseries\songti}}
\makeatother

\geometry{left=2.5cm,right=2.5cm,top=3.5cm,bottom=2.5cm}
\setlength{\headheight}{24.64995pt}

\pagestyle{fancy}
\fancyhf{}
\fancyhead[C]{\textbf{实验4 基本图形变换}}
\fancyhead[R]{\thepage}
\renewcommand{\headrulewidth}{2pt}
\renewcommand{\headrule}{\hbox to\headwidth{\leaders\hrule height 2pt\hfill}}

\lstset{
    language=C++,
    basicstyle=\ttfamily\small,
    keywordstyle=\color{blue},
    commentstyle=\color{green!50!black},
    stringstyle=\color{red},
    numbers=left,
    numberstyle=\tiny\color{gray},
    stepnumber=1,
    numbersep=5pt,
    backgroundcolor=\color{gray!10},
    showspaces=false,
    showstringspaces=false,
    showtabs=false,
    frame=single,
    tabsize=2,
    captionpos=b,
    breaklines=true,
    breakatwhitespace=false,
}

\title{\textbf{实验4 基本图形变换}}
\author{程泽凯 \quad 2023112621}
\date{\today}

\begin{document}

\begin{titlepage}
    \thispagestyle{empty}
    \begin{spacing}{1.2}
    \centering
    \vspace*{10mm}
    {\zihao{4} 《计算机图形学及人机交互》\par}
    \vspace{20mm}
    {\bfseries\zihao{1} 实 验 报 告\par}
    \vspace{40mm}
    \songti\zihao{-3}
    \begin{flushleft}
    \setlength{\tabcolsep}{1em}
    \renewcommand\arraystretch{2}
    \begin{tabular}{@{}ll}
        学\quad 号： & \fillin{2023112621} \\
        姓\quad 名： & \fillin{程泽凯} \\
        年\quad 级： & \fillin{2023 级} \\
        专\quad 业： & \fillin{计算机科学与技术} \\
        院\quad 系： & \fillin{计算机与人工智能学院} \\
    \end{tabular}
    \end{flushleft}
    \vfill
    {\zihao{-3} 2025~2026 第 1 学期\par}
    \end{spacing}
\end{titlepage}

\clearpage
\tableofcontents
\clearpage

\songti\zihao{-4}
\onehalfspacing

\section{实验目的、任务及要求}

\subsection{实验目的}

\begin{tcolorbox}
通过本次实验，学习并实现计算机图形学中的基本图形变换和裁剪算法，掌握二维图形的几何变换原理和裁剪技术。
\begin{enumerate}
    \item 掌握二维图形的基本几何变换：平移、旋转、缩放
    \item 理解齐次坐标和变换矩阵的概念
    \item 掌握Liang-Barsky直线裁剪算法的原理与实现
    \item 掌握Sutherland-Hodgman多边形裁剪算法的原理与实现
    \item 了解Weiler-Atherton多边形裁剪算法的基本思想
    \item 学习任意多边形的绘制方法
\end{enumerate}
\end{tcolorbox}

\subsection{实验任务与要求}

\begin{tcolorbox}
\begin{enumerate}
    \item \textbf{任意多边形绘制}：实现交互式绘制任意顶点数的多边形
    \item \textbf{图形平移}：实现图形的平移变换，支持鼠标拖动操作
    \item \textbf{图形旋转}：实现图形绑定任意点的旋转变换，支持自定义旋转中心
    \item \textbf{图形缩放}：实现图形绑定任意点的缩放变换，支持自定义缩放中心
    \item \textbf{直线裁剪}：采用Liang-Barsky算法实现直线段的裁剪
    \item \textbf{多边形裁剪}：采用Sutherland-Hodgman算法和Weiler-Atherton算法实现多边形裁剪
\end{enumerate}
\end{tcolorbox}

\clearpage

\section{实验环境}
\begin{tcolorbox}
\subsection{硬件环境}

\begin{enumerate}
    \item 计算机：机械革命旷世 16pro
    \item CPU：12th Gen Intel® i9-12900H
    \item RAM：32GB
    \item 显卡：NVIDIA GeForce RTX 4070 Laptop GPU
\end{enumerate}

\subsection{软件环境}

\begin{enumerate}
    \item 操作系统：Windows 11 专业版
    \item 开发工具：Visual Studio 2022
    \item 编程语言：C++17
    \item 图形API：Direct2D (Windows SDK)
    \item 项目类型：Win32桌面应用程序
\end{enumerate}
\end{tcolorbox}

\clearpage

\section{主要功能模块结构及说明}
\begin{tcolorbox}
\subsection{系统总体架构}

本系统基于实验3的绘图软件扩展，采用模块化设计。根据实验任务要求，系统主要包括多边形绘制、图形变换（平移、旋转、缩放）和图形裁剪（直线裁剪、多边形裁剪）三大功能模块。

\begin{figure}[H]
\centering
\begin{tikzpicture}[
    every node/.style={font=\scriptsize},
    mainbox/.style={rectangle, draw=black, line width=1.2pt, fill=yellow!25, minimum width=3cm, minimum height=0.8cm, text centered, rounded corners=3pt, font=\footnotesize\bfseries},
    modulebox/.style={rectangle, draw=black, line width=1pt, fill=blue!20, minimum width=1.6cm, minimum height=0.7cm, text centered, rounded corners=2pt, align=center, font=\scriptsize},
    subbox/.style={rectangle, draw=black, line width=0.8pt, fill=green!15, minimum width=1.4cm, minimum height=0.55cm, text centered, rounded corners=2pt, font=\tiny}
]

% 第一层：主程序
\node[mainbox] (main) at (0, 0) {图形变换系统};

% 第二层：三个功能模块
\node[modulebox] (m1) at (-4.5, -1.8) {多边形绘制};
\node[modulebox] (m2) at (0, -1.8) {图形变换};
\node[modulebox] (m3) at (4.5, -1.8) {图形裁剪};

% 第三层：子模块
\node[subbox] (s1) at (-4.5, -3.5) {任意多边形};
\node[subbox] (s2a) at (-1.5, -3.5) {平移变换};
\node[subbox] (s2b) at (0, -3.5) {旋转变换};
\node[subbox] (s2c) at (1.5, -3.5) {缩放变换};
\node[subbox] (s3a) at (3.3, -3.5) {Liang-Barsky};
\node[subbox] (s3b) at (5, -3.5) {S-H裁剪};
\node[subbox] (s3c) at (6.5, -3.5) {W-A裁剪};

% 连线：主程序到三个模块
\draw[->, line width=0.8pt] ([xshift=-1cm]main.south) -- (m1.north);
\draw[->, line width=0.8pt] (main.south) -- (m2.north);
\draw[->, line width=0.8pt] ([xshift=1cm]main.south) -- (m3.north);

% 连线：模块到子模块
\draw[->, line width=0.6pt] (m1.south) -- (s1.north);
\draw[->, line width=0.6pt] ([xshift=-0.4cm]m2.south) -- (s2a.north);
\draw[->, line width=0.6pt] (m2.south) -- (s2b.north);
\draw[->, line width=0.6pt] ([xshift=0.4cm]m2.south) -- (s2c.north);
\draw[->, line width=0.6pt] ([xshift=-0.4cm]m3.south) -- (s3a.north);
\draw[->, line width=0.6pt] (m3.south) -- (s3b.north);
\draw[->, line width=0.6pt] ([xshift=0.4cm]m3.south) -- (s3c.north);

\end{tikzpicture}
\caption{系统模块结构图}
\label{fig:system_arch}
\end{figure}
\end{tcolorbox}
\subsection{核心模块详细说明}


\begin{tcolorbox}
\subsubsection{多边形绘制模块}
\textbf{模块功能}：

本模块实现任意多边形的交互式绘制功能，用户可以通过鼠标点击依次添加顶点，构建任意形状的多边形。

\textbf{主要功能}：
\begin{itemize}
    \item \textbf{顶点添加}：鼠标左键点击添加多边形顶点
    \item \textbf{实时预览}：绘制过程中实时显示多边形轮廓
    \item \textbf{自动闭合}：右键完成绘制时自动连接首尾顶点
    \item \textbf{填充支持}：支持对封闭多边形进行颜色填充
\end{itemize}

\textbf{用户交互}：
\begin{itemize}
    \item 选择"绘图工具→自由多边形"菜单
    \item 在画布上依次点击添加顶点
    \item 右键完成绘制，多边形自动闭合
\end{itemize}
\end{tcolorbox}
\begin{tcolorbox}
\subsubsection{图形变换模块}

\textbf{模块功能}：

本模块实现二维图形的基本几何变换，包括平移、旋转和缩放三种变换操作。所有变换都基于齐次坐标和变换矩阵实现。

\textbf{平移变换}：
\begin{itemize}
    \item 将图形沿指定方向移动指定距离
    \item 变换矩阵：$T(t_x, t_y) = \begin{pmatrix} 1 & 0 & t_x \\ 0 & 1 & t_y \\ 0 & 0 & 1 \end{pmatrix}$
    \item 用户通过鼠标拖动实现平移操作
\end{itemize}

\textbf{旋转变换}：
\begin{itemize}
    \item 将图形绑定任意点进行旋转
    \item 变换矩阵：$R(\theta) = \begin{pmatrix} \cos\theta & -\sin\theta & 0 \\ \sin\theta & \cos\theta & 0 \\ 0 & 0 & 1 \end{pmatrix}$
    \item 支持自定义旋转中心点
    \item 用户通过鼠标滚轮调整旋转角度
\end{itemize}

\textbf{缩放变换}：
\begin{itemize}
    \item 将图形绑定任意点进行缩放
    \item 变换矩阵：$S(s_x, s_y) = \begin{pmatrix} s_x & 0 & 0 \\ 0 & s_y & 0 \\ 0 & 0 & 1 \end{pmatrix}$
    \item 支持自定义缩放中心点
    \item 用户通过鼠标滚轮调整缩放比例
\end{itemize}

\textbf{用户交互}：
\begin{itemize}
    \item 先选中要变换的图形
    \item 选择"变换"菜单中的平移/旋转/缩放
    \item 平移：拖动鼠标移动图形
    \item 旋转/缩放：可先点击设置中心点，然后滚动滚轮调整
    \item 左键确认变换，右键取消
\end{itemize}
\end{tcolorbox}

\begin{tcolorbox}	\subsubsection{图形裁剪模块}
\textbf{模块功能}：

本模块实现图形的裁剪功能，包括直线裁剪和多边形裁剪两类算法。

\textbf{Liang-Barsky直线裁剪}：
\begin{itemize}
    \item 基于参数化直线方程的裁剪算法
    \item 通过计算直线与窗口边界的交点参数实现裁剪
    \item 效率高于Cohen-Sutherland算法
    \item 适用于矩形裁剪窗口
\end{itemize}

\textbf{Sutherland-Hodgman多边形裁剪}：
\begin{itemize}
    \item 逐边裁剪算法，依次对四条边界进行裁剪
    \item 适用于凸多边形裁剪窗口
    \item 可处理凸多边形和凹多边形
    \item 输出单个连续的多边形
\end{itemize}

\textbf{Weiler-Atherton多边形裁剪}：
\begin{itemize}
    \item 基于交点链表的裁剪算法
    \item 可处理凹多边形和自相交多边形
    \item 支持带孔洞的多边形
    \item 可能输出多个独立的多边形
\end{itemize}

\textbf{用户交互}：
\begin{itemize}
    \item 选择"裁剪"菜单中的裁剪算法
    \item 点击选择要裁剪的图形
    \item 点击选择裁剪矩形
    \item 系统自动执行裁剪并显示结果
\end{itemize}
\end{tcolorbox}

\clearpage

\section{实验算法及程序流程}
\begin{tcolorbox}
\subsection{图形变换算法}

\subsubsection{平移变换}

\textbf{输入}：图形顶点集合$\{P_i\}$，平移向量$(t_x, t_y)$

\textbf{输出}：平移后的顶点集合$\{P'_i\}$

\textbf{关键变量}：$n$=顶点总数，$i$=当前索引，$(t_x, t_y)$=平移向量

\textbf{详细算法流程}：
\begin{enumerate}
    \item \textbf{初始化}：$i = 0$，$n = $顶点总数
    \item \textbf{循环}（当$i < n$）：读取$(x_i, y_i)$，计算$x'_i = x_i + t_x$，$y'_i = y_i + t_y$，$i = i + 1$
    \item \textbf{输出}：返回$\{P'_i\}$
\end{enumerate}

\textbf{变换公式}：
$$\begin{pmatrix} x' \\ y' \\ 1 \end{pmatrix} = \begin{pmatrix} 1 & 0 & t_x \\ 0 & 1 & t_y \\ 0 & 0 & 1 \end{pmatrix} \begin{pmatrix} x \\ y \\ 1 \end{pmatrix}$$

\begin{figure}[H]
\centering
\begin{tikzpicture}[node distance=1.2cm,
    startstop/.style={rectangle, rounded corners, minimum width=2.5cm, minimum height=0.8cm, text centered, draw=black, fill=red!30, font=\small},
    process/.style={rectangle, minimum width=2.8cm, minimum height=0.8cm, text centered, draw=black, fill=orange!30, font=\small, align=center},
    decision/.style={diamond, aspect=2, minimum width=2.5cm, minimum height=0.8cm, text centered, draw=black, fill=green!30, font=\small, align=center},
    arrow/.style={thick,->,>=stealth}
]
\node (start) [startstop] {开始};
\node (init) [process, below=of start] {输入: 顶点$\{P_i\}$, 平移量$(t_x,t_y)$\\初始化: $i=0$, $n=$顶点总数};
\node (loop) [process, below=of init] {取当前顶点$P_i=(x_i,y_i)$};
\node (calc) [process, below=of loop] {计算新坐标:\\$x'_i = x_i + t_x$\\$y'_i = y_i + t_y$\\$i = i + 1$};
\node (check) [decision, below=of calc] {$i < n$?};
\node (result) [process, right=4cm of check] {输出: 新顶点集合$\{P'_i\}$};
\node (end) [startstop, below=of result] {结束};

\draw [arrow] (start) -- (init);
\draw [arrow] (init) -- (loop);
\draw [arrow] (loop) -- (calc);
\draw [arrow] (calc) -- (check);
\draw [arrow] (check) -- node[above] {否} (result);
\draw [arrow] (check.west) -- node[left,yshift=10pt] {是} ++(-3,0) |- (loop);
\draw [arrow] (result) -- (end);
\end{tikzpicture}
\caption{平移变换流程图}
\label{fig:translate_flow}
\end{figure}

\subsubsection{旋转变换}

\textbf{输入}：图形顶点集合$\{P_i\}$，旋转中心$(c_x, c_y)$，旋转角度$\theta$

\textbf{输出}：旋转后的顶点集合$\{P'_i\}$

\textbf{关键变量}：$n$=顶点数，$(c_x, c_y)$=旋转中心，$\theta$=旋转角度，$s=\sin\theta$，$c=\cos\theta$

\textbf{详细算法流程}：
\begin{enumerate}
    \item \textbf{初始化}：$i = 0$，$n = $顶点数，$s = \sin\theta$，$c = \cos\theta$
    \item \textbf{循环}（当$i < n$）：
    \begin{enumerate}
        \item 平移到原点：$x' = x_i - c_x$，$y' = y_i - c_y$
        \item 绕原点旋转：$x'' = x' \cdot c - y' \cdot s$，$y'' = x' \cdot s + y' \cdot c$
        \item 平移回原位：$x'_i = x'' + c_x$，$y'_i = y'' + c_y$
        \item $i = i + 1$
    \end{enumerate}
    \item \textbf{输出}：返回$\{P'_i\}$
\end{enumerate}

\textbf{复合变换矩阵}：
$$T = T(c_x, c_y) \cdot R(\theta) \cdot T(-c_x, -c_y)$$

\begin{figure}[H]
\centering
\begin{tikzpicture}[node distance=1.2cm,
    startstop/.style={rectangle, rounded corners, minimum width=2.5cm, minimum height=0.8cm, text centered, draw=black, fill=red!30, font=\small},
    process/.style={rectangle, minimum width=3cm, minimum height=0.8cm, text centered, draw=black, fill=orange!30, font=\small, align=center},
    decision/.style={diamond, aspect=2, minimum width=2.5cm, minimum height=0.8cm, text centered, draw=black, fill=green!30, font=\small, align=center},
    arrow/.style={thick,->,>=stealth}
]
\node (start) [startstop] {开始};
\node (init) [process, below=of start] {输入: 顶点$\{P_i\}$, 中心$(c_x,c_y)$, 角度$\theta$\\初始化: $i=0$, $n=$顶点数, $s=\sin\theta$, $c=\cos\theta$};
\node (loop) [process, below=of init] {取当前顶点$P_i=(x_i,y_i)$};
\node (trans1) [process, below=of loop] {平移到原点:\\$x' = x_i - c_x$\\$y' = y_i - c_y$};
\node (rotate) [process, below=of trans1] {绕原点旋转:\\$x'' = x' \cdot c - y' \cdot s$\\$y'' = x' \cdot s + y' \cdot c$};
\node (trans2) [process, below=of rotate] {平移回原位置:\\$x'_i = x'' + c_x$, $y'_i = y'' + c_y$\\$i = i + 1$};
\node (check) [decision, below=of trans2] {$i < n$?};
\node (end) [startstop, right=4cm of check] {结束, 输出$\{P'_i\}$};

\draw [arrow] (start) -- (init);
\draw [arrow] (init) -- (loop);
\draw [arrow] (loop) -- (trans1);
\draw [arrow] (trans1) -- (rotate);
\draw [arrow] (rotate) -- (trans2);
\draw [arrow] (trans2) -- (check);
\draw [arrow] (check) -- node[above] {否} (end);
\draw [arrow] (check.west) -- node[left,yshift=10pt] {是} ++(-3,0) |- (loop);
\end{tikzpicture}
\caption{旋转变换流程图}
\label{fig:rotate_flow}
\end{figure}

\subsubsection{缩放变换}

\textbf{输入}：图形顶点集合$\{P_i\}$，缩放中心$(c_x, c_y)$，缩放因子$s$

\textbf{输出}：缩放后的顶点集合$\{P'_i\}$

\textbf{关键变量}：$n$=顶点数，$(c_x, c_y)$=缩放中心，$s$=缩放因子

\textbf{详细算法流程}：
\begin{enumerate}
    \item \textbf{初始化}：$i = 0$，$n = $顶点数
    \item \textbf{循环}（当$i < n$）：
    \begin{enumerate}
        \item 计算相对偏移：$dx = x_i - c_x$，$dy = y_i - c_y$
        \item 缩放偏移：$dx' = dx \times s$，$dy' = dy \times s$
        \item 计算新坐标：$x'_i = c_x + dx'$，$y'_i = c_y + dy'$
        \item $i = i + 1$
    \end{enumerate}
    \item \textbf{输出}：返回$\{P'_i\}$
\end{enumerate}

\textbf{复合变换矩阵}：
$$T = T(c_x, c_y) \cdot S(s, s) \cdot T(-c_x, -c_y)$$

\begin{figure}[H]
\centering
\begin{tikzpicture}[node distance=1.2cm,
    startstop/.style={rectangle, rounded corners, minimum width=2.5cm, minimum height=0.8cm, text centered, draw=black, fill=red!30, font=\small},
    process/.style={rectangle, minimum width=3cm, minimum height=0.8cm, text centered, draw=black, fill=orange!30, font=\small, align=center},
    decision/.style={diamond, aspect=2, minimum width=2.5cm, minimum height=0.8cm, text centered, draw=black, fill=green!30, font=\small, align=center},
    arrow/.style={thick,->,>=stealth}
]
\node (start) [startstop] {开始};
\node (init) [process, below=of start] {输入: 顶点$\{P_i\}$, 中心$(c_x,c_y)$, 因子$s$\\初始化: $i=0$, $n=$顶点数};
\node (loop) [process, below=of init] {取当前顶点$P_i=(x_i,y_i)$};
\node (offset) [process, below=of loop] {计算相对偏移:\\$dx = x_i - c_x$\\$dy = y_i - c_y$};
\node (scale) [process, below=of offset] {缩放偏移量:\\$dx' = dx \times s$\\$dy' = dy \times s$};
\node (newpos) [process, below=of scale] {计算新坐标:\\$x'_i = c_x + dx'$, $y'_i = c_y + dy'$\\$i = i + 1$};
\node (check) [decision, below=of newpos] {$i < n$?};
\node (end) [startstop, right=4cm of check] {结束, 输出$\{P'_i\}$};

\draw [arrow] (start) -- (init);
\draw [arrow] (init) -- (loop);
\draw [arrow] (loop) -- (offset);
\draw [arrow] (offset) -- (scale);
\draw [arrow] (scale) -- (newpos);
\draw [arrow] (newpos) -- (check);
\draw [arrow] (check) -- node[above] {否} (end);
\draw [arrow] (check.west) -- node[left,yshift=10pt] {是} ++(-3,0) |- (loop);
\end{tikzpicture}
\caption{缩放变换流程图}
\label{fig:scale_flow}
\end{figure}

\clearpage
\end{tcolorbox}

\begin{tcolorbox}
\subsection{Liang-Barsky直线裁剪算法}

\textbf{输入}：直线端点$(x_1, y_1)$和$(x_2, y_2)$，裁剪窗口$(x_{min}, y_{min}, x_{max}, y_{max})$

\textbf{输出}：裁剪后的直线端点，或空（直线完全在窗口外）

\textbf{算法原理}：

直线的参数方程为：
$$x = x_1 + t \cdot \Delta x, \quad y = y_1 + t \cdot \Delta y, \quad t \in [0, 1]$$

其中$\Delta x = x_2 - x_1$，$\Delta y = y_2 - y_1$。

对于四条边界，定义参数$p_k$和$q_k$：
\begin{align*}
p_1 &= -\Delta x, \quad q_1 = x_1 - x_{min} \quad \text{(左边界)} \\
p_2 &= \Delta x, \quad q_2 = x_{max} - x_1 \quad \text{(右边界)} \\
p_3 &= -\Delta y, \quad q_3 = y_1 - y_{min} \quad \text{(下边界)} \\
p_4 &= \Delta y, \quad q_4 = y_{max} - y_1 \quad \text{(上边界)}
\end{align*}

\textbf{关键变量}：
\begin{itemize}
    \item $\Delta x = x_2 - x_1$，$\Delta y = y_2 - y_1$：方向向量
    \item $t_{enter}$：进入参数（初始为0），$t_{leave}$：离开参数（初始为1）
    \item $p_k, q_k$：第$k$条边界的参数（$k=1,2,3,4$对应左/右/下/上）
    \item $t$：当前边界的交点参数
\end{itemize}

\textbf{详细算法流程}：
\begin{enumerate}
    \item \textbf{初始化}：$\Delta x = x_2 - x_1$，$\Delta y = y_2 - y_1$，$t_{enter} = 0$，$t_{leave} = 1$，$k = 1$
    \item \textbf{循环}（当$k \leq 4$）：
    \begin{enumerate}
        \item 计算$p_k$和$q_k$（根据边界类型）
        \item 若$p_k = 0$且$q_k < 0$：直线平行于边界且在外侧，返回空
        \item 若$p_k \neq 0$：计算$t = q_k / p_k$
        \item 若$p_k < 0$（从外到内）：$t_{enter} = \max(t_{enter}, t)$
        \item 若$p_k > 0$（从内到外）：$t_{leave} = \min(t_{leave}, t)$
        \item $k = k + 1$
    \end{enumerate}
    \item \textbf{判断}：若$t_{enter} > t_{leave}$，返回空（无交集）
    \item \textbf{计算裁剪端点}：$(x'_1, y'_1) = (x_1 + t_{enter} \cdot \Delta x, y_1 + t_{enter} \cdot \Delta y)$，$(x'_2, y'_2) = (x_1 + t_{leave} \cdot \Delta x, y_1 + t_{leave} \cdot \Delta y)$
\end{enumerate}

\begin{figure}[H]
	\centering
	\begin{tikzpicture}[
		node distance=0.9cm, % 微调间距，平衡紧凑度和可读性
		startstop/.style={rectangle, rounded corners, minimum width=2cm, minimum height=0.7cm, text centered, draw=black, fill=red!30, font=\small},
		process/.style={rectangle, minimum width=2.8cm, minimum height=0.7cm, text centered, draw=black, fill=orange!30, font=\small, align=center},
		decision/.style={diamond, aspect=2, minimum width=2.2cm, minimum height=0.7cm, text centered, draw=black, fill=green!30, font=\small, align=center},
		arrow/.style={thick,->,>=stealth}
		]
		% 顶层节点
		\node (start) [startstop] {开始};
		\node (init) [process, below=of start] {输入: 端点$(x_1,y_1),(x_2,y_2)$, 窗口边界\\初始化: $\Delta x=x_2-x_1$, $\Delta y=y_2-y_1$\\$t_{enter}=0$, $t_{leave}=1$, $k=1$};
		\node (loop) [process, below=of init] {计算第$k$条边界参数:\\$p_k$, $q_k$ (左/右/下/上)};
		
		% 条件判断节点（垂直排列，保持对齐）
		\node (checkp) [decision, below=of loop] {$p_k = 0$?};
		\node (checkq) [decision, below=of checkp] {$q_k < 0$?}; % 取消额外间距，避免纵向拉伸
		\node (reject1) [process, right=2cm of checkq] {返回: 直线在窗口外};
		\node (calct) [process, below=1cm of checkq] {计算交点参数:\\$t = q_k / p_k$}; % 与checkq保持合理间距
		\node (checkdir) [decision, below=of calct] {$p_k < 0$?};
		
		% 输入/输出更新节点（左右紧凑对齐）
		\node (enter) [process, left=2cm of checkdir] {更新进入点:\\$t_{enter} = \max(t_{enter}, t)$};
		\node (leave) [process, right=2cm of checkdir] {更新离开点:\\$t_{leave} = \min(t_{leave}, t)$};
		
		% 最终判断+输出节点
		\node (checkfinal) [decision, below=1.2cm of checkdir] {$t_{enter} > t_{leave}$?};
		\node (reject2) [process, right=2cm of checkfinal] {返回: 直线在窗口外};
		\node (calc) [process, below=of checkfinal] {计算裁剪后端点:\\$(x'_1,y'_1)=(x_1+t_{enter}\Delta x, y_1+t_{enter}\Delta y)$\\$(x'_2,y'_2)=(x_1+t_{leave}\Delta x, y_1+t_{leave}\Delta y)$};
		\node (end) [startstop, below=of calc] {结束, 输出裁剪后端点};
		
		% 关键修改：调整checkp的“否”分支路径，绕开checkq节点
		\draw [arrow] (start) -- (init);
		\draw [arrow] (init) -- (loop);
		\draw [arrow] (loop) -- (checkp);
		
		% p_k=0? 的“是”分支（垂直向下连checkq）
		\draw [arrow] (checkp) -- node[right, pos=0.4] {是} (checkq);
		
		% p_k=0? 的“否”分支（先向左横向延伸，再垂直向下绕开checkq，连calct）
		\draw [arrow] (checkp.west) -- node[above, pos=0.5] {否} ++(-1.8, 0) |- (calct.west);
		
		% 其他连线保持清晰，无重叠
		\draw [arrow] (checkq) -- node[above, pos=0.3] {是} (reject1);
		\draw [arrow] (checkq) -- node[right, pos=0.4] {否} (calct);
		\draw [arrow] (calct) -- (checkdir);
		\draw [arrow] (checkdir) -- node[above, pos=0.3] {是} (enter);
		\draw [arrow] (checkdir) -- node[above, pos=0.3] {否} (leave);
		\draw [arrow] (enter.south) |- (checkfinal.west); % 从enter右侧折回，避免左偏
\draw [arrow] (leave.south) -- ++(0, -0.2) % 向下短距离延伸（与enter对称）
-- (checkfinal.north) % 水平向左直接连checkfinal右上角
-- (checkfinal.north);
		\draw [arrow] (checkfinal) -- node[above, pos=0.3] {是} (reject2);
		\draw [arrow] (checkfinal) -- node[right, pos=0.4] {否} (calc);
		\draw [arrow] (calc) -- (end);
	\end{tikzpicture}
	\caption{优化后的Liang-Barsky直线裁剪算法流程图}
	\label{fig:liang_barsky_flow_optimized}
\end{figure}

\end{tcolorbox}

\clearpage

\begin{tcolorbox}
\subsection{Sutherland-Hodgman多边形裁剪算法}

\textbf{输入}：多边形顶点列表$\{P_0, P_1, \ldots, P_{n-1}\}$，裁剪窗口$(x_{min}, y_{min}, x_{max}, y_{max})$

\textbf{输出}：裁剪后的多边形顶点列表

\textbf{算法原理}：

逐边裁剪算法，依次对裁剪窗口的四条边界进行裁剪。对于每条边界，遍历多边形的每条边，根据边的两个端点与边界的位置关系，决定输出哪些顶点。

\textbf{关键变量}：
\begin{itemize}
    \item $output$：当前输出多边形顶点列表
    \item $edge$：当前处理的边界索引（0=下，1=右，2=上，3=左）
    \item $P_i, P_{i+1}$：当前处理的边的起点和终点
    \item $I$：边与边界的交点
\end{itemize}

\textbf{四种情况}：
\begin{enumerate}
    \item 两点都在内侧：输出$P_{i+1}$
    \item $P_i$内、$P_{i+1}$外：计算交点$I$，输出$I$
    \item 两点都在外侧：不输出
    \item $P_i$外、$P_{i+1}$内：计算交点$I$，输出$I$和$P_{i+1}$
\end{enumerate}

\textbf{详细算法流程}：
\begin{enumerate}
    \item \textbf{初始化}：$output = \{P_i\}$（输入多边形），$edge = 0$
    \item \textbf{外循环}（当$edge < 4$）：
    \begin{enumerate}
        \item 选择第$edge$条边界
        \item \textbf{内循环}：对$output$中每条边$(P_i, P_{i+1})$：
        \begin{enumerate}
            \item 判断$P_i$和$P_{i+1}$与边界的位置关系
            \item 根据四种情况决定输出顶点
        \end{enumerate}
        \item 更新$output$为新的顶点列表
        \item $edge = edge + 1$
    \end{enumerate}
    \item \textbf{输出}：返回$output$
\end{enumerate}

\begin{figure}[H]
	\centering
	\begin{tikzpicture}[
		node distance=0.9cm, % 缩小节点间距
		startstop/.style={rectangle, rounded corners, minimum width=2cm, minimum height=0.7cm, text centered, draw=black, fill=red!30, font=\small},
		process/.style={rectangle, minimum width=2.8cm, minimum height=0.7cm, text centered, draw=black, fill=orange!30, font=\small, align=center},
		decision/.style={diamond, aspect=2, minimum width=2.2cm, minimum height=0.7cm, text centered, draw=black, fill=green!30, font=\small, align=center},
		arrow/.style={thick,->,>=stealth}
		]
		% 顶层节点
		\node (start) [startstop] {开始};
		\node (init) [process, below=of start] {输入: 多边形$\{P_i\}$, 窗口边界\\初始化: $output=\{P_i\}$, $edge=0$};
		\node (edge) [process, below=of init] {选择第$edge$条边界(下/右/上/左)\\$edge = edge + 1$};
		
		% 循环主体（取边+位置判断）
		\node (vertex) [process, below=of edge] {取边$(P_i, P_{i+1})$\\判断与边界的位置关系};
		
		% 条件判断节点（垂直紧凑排列）
		\node (case1) [decision, below=of vertex] {两点都在内侧?};
		\node (case2) [decision, below=1.2cm of case1] {$P_i$内、$P_{i+1}$外?};
		\node (case3) [decision, below=1.2cm of case2] {$P_i$外、$P_{i+1}$内?};
		\node (next) [decision, below=1.2cm of case3] {$i < n-1$?};
		\node (nextedge) [decision, below=1.2cm of next] {$edge < 4$?};
		
		% 输出操作节点（右对齐，与条件节点对应）
		\node (out1) [process, right=2.5cm of case1] {将$P_{i+1}$加入输出};
		\node (out2) [process, right=2.5cm of case2] {计算交点$I$并加入输出};
		\node (out3) [process, right=2.5cm of case3] {计算交点$I$, 将$I$和$P_{i+1}$加入输出};
		
		% 结束节点
		\node (end) [startstop, right=2.5cm of nextedge] {结束, 输出裁剪后多边形};
		
		
		% 连线（简化交叉，用垂直/水平折线）
		\draw [arrow] (start) -- (init);
		\draw [arrow] (init) -- (edge);
		\draw [arrow] (edge) -- (vertex);
		\draw [arrow] (vertex) -- (case1);
		
		% 条件1分支
		\draw [arrow] (case1) -- node[above, pos=0.3] {是} (out1);
		\draw [arrow] (case1) -- node[right, pos=0.2] {否} (case2);
		
		% 条件2分支
		\draw [arrow] (case2) -- node[above, pos=0.3] {是} (out2);
		\draw [arrow] (case2) -- node[right, pos=0.2] {否} (case3);
		
		% 条件3分支
		\draw [arrow] (case3) -- node[above, pos=0.3] {是} (out3);
		\draw [arrow] (case3) -- node[right, pos=0.2] {否} (next);
		
		% 输出操作→i循环判断
		\draw [arrow] (out1) |- (next);
		\draw [arrow] (out2) |- (next);
		\draw [arrow] (out3) |- (next);
		
		% i循环分支（回到“取边”）
		\draw [arrow] (next.west) -- ++(-1.5,0) node[above, pos=0.5] {是} |- (vertex.west);
		\draw [arrow] (next) -- node[right, pos=0.2] {否} (nextedge);
		
		% edge循环分支（回到“选边界”）
		\draw [arrow] (nextedge.west) -- ++(-2,0) node[above, pos=0.5] {是} |- (edge.west);
		\draw [arrow] (nextedge) -- node[above, pos=0.3] {否} (end);
	\end{tikzpicture}
	\caption{Sutherland-Hodgman多边形裁剪算法流程图}
	\label{fig:sh_flow_optimized}
\end{figure}

\end{tcolorbox}

\clearpage

\begin{tcolorbox}
\subsection{Weiler-Atherton多边形裁剪算法}

\textbf{输入}：主多边形（被裁剪多边形）顶点列表$S=\{S_0, S_1, ..., S_{n-1}\}$，裁剪多边形顶点列表$C=\{C_0, C_1, ..., C_{m-1}\}$

\textbf{输出}：裁剪后的多边形列表$R$（凹多边形可能产生多个独立区域）

\textbf{关键变量说明}：
\begin{itemize}
    \item $subjList$：主多边形顶点链表，包含原始顶点和插入的交点
    \item $clipList$：裁剪多边形顶点链表，包含原始顶点和插入的交点
    \item $intersections$：所有交点信息列表，每个交点包含：
    \begin{itemize}
        \item $pt$：交点坐标$(x, y)$
        \item $t_S$：交点在主多边形边上的参数位置$t_S \in (0,1)$
        \item $t_C$：交点在裁剪多边形边上的参数位置$t_C \in (0,1)$
        \item $entering$：是否为进入点（true/false）
    \end{itemize}
    \item $neighbor$：交点在另一个链表中对应节点的索引（用于切换链表）
    \item $visited$：标记交点是否已被访问（避免重复遍历）
\end{itemize}

\textbf{算法原理}：

Weiler-Atherton算法基于交点链表，可以处理凹多边形、自相交多边形和带孔洞的多边形。

\textbf{交点分类}：
\begin{itemize}
    \item \textbf{进入点(Entering)}：主多边形从裁剪多边形外部进入内部的交点
    \item \textbf{离开点(Leaving)}：主多边形从裁剪多边形内部离开到外部的交点
\end{itemize}

\textbf{详细算法流程}：

\textbf{步骤1：计算所有交点}
\begin{enumerate}
    \item 初始化：$intersections = \emptyset$
    \item 对于主多边形的每条边$e_S = (S_i, S_{(i+1)\%n})$，$i = 0, 1, ..., n-1$：
    \begin{enumerate}
        \item 对于裁剪多边形的每条边$e_C = (C_j, C_{(j+1)\%m})$，$j = 0, 1, ..., m-1$：
        \item 计算$e_S$与$e_C$的交点，使用参数方程：$P = S_i + t_S \cdot (S_{i+1} - S_i)$
        \item 若存在有效交点（$0 < t_S < 1$ 且 $0 < t_C < 1$）：
        \item \quad 计算交点坐标：$pt = S_i + t_S \cdot (S_{i+1} - S_i)$
        \item \quad 判断进入/离开：取交点后一小段位置$pt' = S_i + (t_S + \epsilon) \cdot (S_{i+1} - S_i)$
        \item \quad 若$pt'$在裁剪多边形内部，则为进入点；否则为离开点
        \item \quad 将交点信息加入$intersections$
    \end{enumerate}
\end{enumerate}

\textbf{步骤2：构建主多边形顶点链表}
\begin{enumerate}
    \item 初始化：$subjList = \emptyset$
    \item 对于每个原始顶点$S_i$，$i = 0, 1, ..., n-1$：
    \begin{enumerate}
        \item 将$S_i$加入$subjList$（标记为非交点）
        \item 收集边$(S_i, S_{i+1})$上的所有交点，按$t_S$升序排序
        \item 依次将这些交点插入$subjList$（标记为交点，设置$entering$属性）
    \end{enumerate}
\end{enumerate}

\textbf{步骤3：构建裁剪多边形顶点链表}（同步骤2，对裁剪多边形执行）

\textbf{步骤4：建立交点双向链接}
\begin{enumerate}
    \item 对于$subjList$中的每个交点$I_S$：
    \begin{enumerate}
        \item 在$clipList$中找到坐标相同的交点$I_C$
        \item 设置$I_S.neighbor = I_C$的索引，$I_C.neighbor = I_S$的索引
    \end{enumerate}
\end{enumerate}

\textbf{步骤5：遍历生成裁剪结果}
\begin{enumerate}
    \item 初始化：$R = \emptyset$（结果多边形列表）
    \item 对于$subjList$中的每个未访问的进入点$I_{start}$：
    \begin{enumerate}
        \item 初始化当前多边形：$polygon = \emptyset$
        \item 设置当前位置：$current = I_{start}$，$onSubject = true$
        \item 循环直到回到起点：
        \begin{enumerate}
            \item 若$onSubject = true$（在主多边形上）：
            \item \quad 将$current.pt$加入$polygon$
            \item \quad 标记$current.visited = true$
            \item \quad 若$current$是离开点且有链接：
            \item \quad \quad 切换到裁剪多边形：$current = clipList[current.neighbor + 1]$
            \item \quad \quad 设置$onSubject = false$
            \item \quad 否则：移动到下一个节点$current = subjList[(current+1) \% |subjList|]$
            \item 若$onSubject = false$（在裁剪多边形上）：
            \item \quad 若$current$是交点且对应主多边形上的进入点：
            \item \quad \quad 切换回主多边形：$current = subjList[current.neighbor]$
            \item \quad \quad 设置$onSubject = true$
            \item \quad 否则：将$current.pt$加入$polygon$，移动到下一个节点
            \item 若回到起点$I_{start}$，退出循环
        \end{enumerate}
        \item 若$|polygon| \geq 3$，将$polygon$加入结果$R$
    \end{enumerate}
    \item 返回$R$
\end{enumerate}

\begin{figure}[H]
	\centering
	\begin{tikzpicture}[node distance=1.2cm,
		startstop/.style={rectangle, rounded corners, minimum width=2.5cm, minimum height=0.8cm, text centered, draw=black, fill=red!30, font=\small},
		process/.style={rectangle, minimum width=3cm, minimum height=0.8cm, text centered, draw=black, fill=orange!30, font=\small, align=center},
		decision/.style={diamond, aspect=2, minimum width=2.5cm, minimum height=0.8cm, text centered, draw=black, fill=green!30, font=\small, align=center},
		arrow/.style={thick,->,>=stealth}
		]
		\node (start) [startstop] {开始};
		\node (init) [process, below=of start] {输入: 主多边形$S$, 裁剪多边形$C$\\初始化: 结果列表$R=\emptyset$};
		\node (intersect) [process, below=of init] {计算$S$与$C$所有边的交点\\记录交点在各边上的参数$t$};
		\node (insert) [process, below=of intersect] {按参数$t$顺序将交点插入\\$S$和$C$的顶点链表};
		\node (mark) [process, below=of insert] {遍历交点, 判断$S$的行进方向\\标记为进入点或离开点};
		\node (find) [decision, below=of mark] {存在未访问的进入点?};
		\node (walk) [process, below=1.5cm of find] {选择一个未访问进入点$I$\\标记$I$为已访问, 沿$S$行走};
		\node (switch) [process, below=of walk] {遇到交点时切换到另一多边形\\继续沿当前多边形方向行走};
		\node (back) [decision, below=of switch] {回到起点$I$?};
		\node (output) [process, right=3cm of back] {将路径上的点构成多边形\\加入结果列表$R$};
		\node (end) [startstop, right=3cm of find] {结束, 输出$R$};
		
		% 原有所有箭头保持不变（除了下面要改的“output→find”）
		\draw [arrow] (start) -- (init);
		\draw [arrow] (init) -- (intersect);
		\draw [arrow] (intersect) -- (insert);
		\draw [arrow] (insert) -- (mark);
		\draw [arrow] (mark) -- (find);
		\draw [arrow] (find) -- node[right] {是} (walk);
		\draw [arrow] (find.east) -- node[above] {否} (end.west);  % 之前改的水平直线
		\draw [arrow] (walk) -- (switch);
		\draw [arrow] (switch) -- (back);
		\draw [arrow] (back) -- node[above] {是} (output);
		\draw [arrow] (back.west) -- node[left,yshift=8pt] {否} ++(-3,0) |- (switch);
		
		% % 输出多边形 → 有未访问的进入点？（折两次：上→左→上）
		\draw [arrow] (output.north) -- ++(0, 6) -- ++(-4.45, 0) -- (find.east);
	\end{tikzpicture}
	\caption{Weiler-Atherton多边形裁剪算法流程图}
	\label{fig:wa_flow}
\end{figure}

\end{tcolorbox}

\clearpage

\section{程序实现}

\begin{tcolorbox}
\subsection{算法设计思路}

\subsubsection{图形变换的设计思路}

图形变换的核心思想是将复杂的变换分解为简单的基本操作。对于绑定任意点的旋转和缩放，采用"三步法"：

\textbf{设计思路}：
\begin{enumerate}
    \item \textbf{问题分析}：直接绕任意点变换较为复杂，但绕原点变换的公式简单明了
    \item \textbf{解决方案}：将任意点变换转化为绕原点变换
    \item \textbf{实现步骤}：
    \begin{itemize}
        \item 第一步：将图形平移，使变换中心移动到原点
        \item 第二步：在原点处执行旋转或缩放
        \item 第三步：将图形平移回原来的位置
    \end{itemize}
\end{enumerate}

这种设计的优点是代码简洁、易于理解，且可以复用基本变换函数。

\subsubsection{裁剪算法的设计思路}

\textbf{Liang-Barsky算法思路}：
\begin{enumerate}
    \item \textbf{核心思想}：将直线表示为参数方程$P(t) = P_1 + t(P_2-P_1)$，$t\in[0,1]$
    \item \textbf{关键观察}：裁剪问题转化为求参数$t$的有效区间
    \item \textbf{算法策略}：
    \begin{itemize}
        \item 对每条边界计算交点参数$t$
        \item 区分"进入"和"离开"两种情况
        \item 进入点取最大值，离开点取最小值
        \item 若进入参数$>$离开参数，则无交集
    \end{itemize}
\end{enumerate}

\textbf{Sutherland-Hodgman算法思路}：
\begin{enumerate}
    \item \textbf{核心思想}：分而治之，逐边裁剪
    \item \textbf{关键观察}：对一条边界裁剪后的结果，可以作为下一条边界裁剪的输入
    \item \textbf{算法策略}：依次对四条边界进行裁剪，每次只考虑一条边界的内外关系
\end{enumerate}

\textbf{Weiler-Atherton算法思路}：
\begin{enumerate}
    \item \textbf{核心思想}：基于交点链表的遍历
    \item \textbf{关键观察}：裁剪结果的边界由主多边形和裁剪多边形的边交替组成
    \item \textbf{算法策略}：
    \begin{itemize}
        \item 计算所有交点，标记为"进入点"或"离开点"
        \item 从进入点开始沿主多边形行走
        \item 遇到离开点时切换到裁剪多边形
        \item 遇到进入点时切换回主多边形
        \item 直到回到起点，形成一个封闭区域
    \end{itemize}
\end{enumerate}

\end{tcolorbox}

\begin{tcolorbox}
\subsection{核心代码实现}

\subsubsection{平移变换实现}

\textbf{代码思路}：

平移是最简单的几何变换，只需将每个顶点的坐标加上平移向量即可。

\textbf{实现要点}：
\begin{itemize}
    \item 平移不改变图形的形状和大小，只改变位置
    \item 对于多边形，需要遍历所有顶点进行平移
    \item 平移后需要同步更新用于碰撞检测的points数组
\end{itemize}

\begin{lstlisting}[caption={平移变换函数}]
// ========== 平移变换 ==========
// 输入: dx, dy - 平移向量
// 输出: 无(直接修改图形顶点)
// 算法: 对每个顶点加上平移向量
void Move(float dx, float dy) override {
    // 遍历所有顶点,加上平移量
    for (size_t i = 0; i < polyPoints.size(); i++) {
        polyPoints[i].x += dx;  // x' = x + dx
        polyPoints[i].y += dy;  // y' = y + dy
    }
    points = polyPoints;  // 同步更新points数组
}
\end{lstlisting}

\subsubsection{旋转变换实现}

\textbf{代码思路}：

绑定任意点的旋转变换采用"三步法"实现：先平移使旋转中心到原点，然后绕原点旋转，最后平移回原位置。

\textbf{实现要点}：
\begin{itemize}
    \item 预先计算$\sin\theta$和$\cos\theta$，避免在循环中重复计算
    \item 绕原点旋转使用标准旋转矩阵：$\begin{pmatrix} \cos\theta & -\sin\theta \\ \sin\theta & \cos\theta \end{pmatrix}$
    \item 注意旋转角度使用弧度制，正值为逆时针旋转
\end{itemize}

\begin{lstlisting}[caption={旋转变换函数}]
// ========== 绕任意点旋转变换 ==========
// 输入: p - 待旋转的点坐标
//       center - 旋转中心点坐标
//       angle - 旋转角度(弧度)
// 输出: 旋转后的点坐标
// 算法: 先平移到原点, 旋转后再平移回去
inline D2D1_POINT_2F RotatePoint(D2D1_POINT_2F p, 
    D2D1_POINT_2F center, float angle) {
    // 步骤1: 预计算三角函数值(避免重复计算)
    float s = sin(angle);  // sin(theta)
    float c = cos(angle);  // cos(theta)
    
    // 步骤2: 平移到原点(相对于旋转中心)
    p.x -= center.x;  // x' = x - cx
    p.y -= center.y;  // y' = y - cy
    
    // 步骤3: 绕原点旋转(应用旋转矩阵)
    // [x'']   [cos -sin] [x']
    // [y''] = [sin  cos] [y']
    float xnew = p.x * c - p.y * s;  // x'' = x'*cos - y'*sin
    float ynew = p.x * s + p.y * c;  // y'' = x'*sin + y'*cos
    
    // 步骤4: 平移回原位置
    return D2D1::Point2F(xnew + center.x, ynew + center.y);
}
\end{lstlisting}

\subsubsection{缩放变换实现}

\textbf{代码思路}：

绑定任意点的缩放变换思路与旋转类似：计算点相对于缩放中心的偏移量，将偏移量乘以缩放因子，再加回中心点坐标。

\textbf{实现要点}：
\begin{itemize}
    \item 缩放因子$>1$表示放大，$<1$表示缩小，$=1$表示不变
    \item 本实现采用等比例缩放（$s_x = s_y$），保持图形比例不变
    \item 代码将三步合并为一行，提高执行效率
\end{itemize}

\begin{lstlisting}[caption={缩放变换函数}]
// ========== 绕任意点缩放变换 ==========
// 输入: p - 待缩放的点坐标
//       center - 缩放中心点坐标
//       factor - 缩放因子(>1放大, <1缩小)
// 输出: 缩放后的点坐标
// 算法: 计算相对偏移, 缩放后加回中心点
inline D2D1_POINT_2F ScalePoint(D2D1_POINT_2F p, 
    D2D1_POINT_2F center, float factor) {
    // 步骤1: 计算相对于中心的偏移量
    // dx = x - cx, dy = y - cy
    // 步骤2: 缩放偏移量
    // dx' = dx * factor, dy' = dy * factor
    // 步骤3: 加回中心点得到新坐标
    // x' = cx + dx', y' = cy + dy'
    return D2D1::Point2F(
        center.x + (p.x - center.x) * factor,  // x' = cx + (x-cx)*s
        center.y + (p.y - center.y) * factor   // y' = cy + (y-cy)*s
    );
}
\end{lstlisting}

\subsubsection{Liang-Barsky直线裁剪实现}

\textbf{代码思路解析}：

Liang-Barsky算法的核心是利用直线的参数方程。设直线从$P_1$到$P_2$，则直线上任意一点可表示为：
$$P(t) = P_1 + t \cdot (P_2 - P_1), \quad t \in [0, 1]$$

对于裁剪窗口的四条边界，我们需要找到直线进入窗口和离开窗口的参数值。关键在于：
\begin{itemize}
    \item 当$p_k < 0$时，直线从边界外侧进入内侧（进入点）
    \item 当$p_k > 0$时，直线从边界内侧离开到外侧（离开点）
    \item 当$p_k = 0$时，直线与边界平行
\end{itemize}

最终，若所有进入点的最大参数值小于所有离开点的最小参数值，则直线与窗口有交集。

\begin{lstlisting}[caption={Liang-Barsky直线裁剪算法}]
// ========== Liang-Barsky直线裁剪算法 ==========
// 输入: (x1,y1),(x2,y2) - 直线端点(会被修改)
//       xmin,ymin,xmax,ymax - 裁剪窗口边界
// 输出: true=有交集(端点已更新), false=无交集
// 原理: 直线参数方程 P(t) = P1 + t*(P2-P1), t in [0,1]
static bool ClipLine(float& x1, float& y1, float& x2, float& y2,
    float xmin, float ymin, float xmax, float ymax) {
    // 步骤1: 计算方向向量
    float dx = x2 - x1;  // delta_x = x2 - x1
    float dy = y2 - y1;  // delta_y = y2 - y1
    
    // 步骤2: 初始化参数区间 [t_enter, t_leave]
    float t_enter = 0.0f;  // 进入窗口的参数(初始为起点)
    float t_leave = 1.0f;  // 离开窗口的参数(初始为终点)
    
    // 步骤3: 定义四条边界的p和q参数
    // p[k]<0: 从外到内, p[k]>0: 从内到外, p[k]=0: 平行
    float p[4] = { -dx, dx, -dy, dy };      // 左/右/下/上
    float q[4] = { x1-xmin, xmax-x1, y1-ymin, ymax-y1 };
    
    // 步骤4: 遍历四条边界,更新参数区间
    for (int i = 0; i < 4; i++) {
        if (p[i] == 0) {
            // 直线平行于该边界
            if (q[i] < 0) return false;  // 在边界外侧,无交集
        } else {
            float t = q[i] / p[i];  // 计算交点参数
            if (p[i] < 0) {
                // 从外到内: 更新进入参数(取较大值)
                t_enter = max(t_enter, t);
            } else {
                // 从内到外: 更新离开参数(取较小值)
                t_leave = min(t_leave, t);
            }
        }
    }
    
    // 步骤5: 判断是否有有效交集
    if (t_enter > t_leave) return false;  // 无交集
    
    // 步骤6: 计算裁剪后的端点坐标
    float new_x1 = x1 + t_enter * dx;  // 新起点
    float new_y1 = y1 + t_enter * dy;
    float new_x2 = x1 + t_leave * dx;  // 新终点
    float new_y2 = y1 + t_leave * dy;
    
    // 步骤7: 更新端点并返回
    x1 = new_x1; y1 = new_y1;
    x2 = new_x2; y2 = new_y2;
    return true;
}
\end{lstlisting}

\end{tcolorbox}

\begin{tcolorbox}
\subsubsection{Sutherland-Hodgman多边形裁剪实现}

\textbf{代码思路}：

Sutherland-Hodgman算法采用"逐边裁剪"的策略，将复杂的多边形裁剪问题分解为四次简单的单边界裁剪。

\textbf{算法核心思想}：
\begin{enumerate}
    \item 依次对下、右、上、左四条边界进行裁剪
    \item 每次裁剪的输出作为下一次裁剪的输入
    \item 对于每条边界，遍历多边形的每条边，根据端点位置决定输出
\end{enumerate}

\textbf{四种情况的处理}：

\begin{table}[H]
\centering
\begin{tabular}{|c|c|c|l|l|}
\hline
\textbf{情况} & \textbf{起点} & \textbf{终点} & \textbf{输出} & \textbf{说明} \\
\hline
1 & 内侧 & 内侧 & 输出终点 & 边完全在内部 \\
2 & 内侧 & 外侧 & 输出交点 & 边从内到外被截断 \\
3 & 外侧 & 外侧 & 不输出 & 边完全在外部 \\
4 & 外侧 & 内侧 & 输出交点和终点 & 边从外到内被截断 \\
\hline
\end{tabular}
\caption{Sutherland-Hodgman四种情况}
\end{table}

\textbf{实现要点}：
\begin{itemize}
    \item 使用取模运算\texttt{(i+1)\%size()}实现多边形的循环遍历
    \item \texttt{IsInside}函数根据边界类型判断点的内外关系
    \item \texttt{ComputeIntersection}函数计算线段与边界的交点
    \item 算法保证输出多边形的连续性和封闭性
\end{itemize}

\begin{lstlisting}[caption={Sutherland-Hodgman多边形裁剪算法}]
// ========== Sutherland-Hodgman多边形裁剪算法 ==========
// 输入: polygon - 待裁剪多边形顶点列表
//       xmin,ymin,xmax,ymax - 裁剪窗口边界
// 输出: 裁剪后的多边形顶点列表
// 原理: 逐边裁剪,依次对四条边界进行处理
static std::vector<D2D1_POINT_2F> ClipPolygon(
    const std::vector<D2D1_POINT_2F>& polygon,
    float xmin, float ymin, float xmax, float ymax) {
    
    if (polygon.size() < 3) return std::vector<D2D1_POINT_2F>();
    
    std::vector<D2D1_POINT_2F> output = polygon;  // 初始化输出
    
    // 依次对四条边界进行裁剪(下->右->上->左)
    output = ClipAgainstEdge(output, xmin, ymin, xmax, ymin, 0);  // 下边界
    output = ClipAgainstEdge(output, xmax, ymin, xmax, ymax, 1);  // 右边界
    output = ClipAgainstEdge(output, xmax, ymax, xmin, ymax, 2);  // 上边界
    output = ClipAgainstEdge(output, xmin, ymax, xmin, ymin, 3);  // 左边界
    
    return output;
}

// ========== 对单条边界进行裁剪 ==========
// 输入: input - 当前多边形, edge - 边界类型(0下/1右/2上/3左)
// 输出: 裁剪后的多边形
// 四种情况: 内内->输出终点, 内外->输出交点
//          外外->不输出, 外内->输出交点和终点
static std::vector<D2D1_POINT_2F> ClipAgainstEdge(
    const std::vector<D2D1_POINT_2F>& input,
    float x1, float y1, float x2, float y2, int edge) {
    
    std::vector<D2D1_POINT_2F> output;
    if (input.empty()) return output;
    
    // 遍历多边形的每条边
    for (size_t i = 0; i < input.size(); i++) {
        D2D1_POINT_2F current = input[i];              // 当前顶点
        D2D1_POINT_2F next = input[(i + 1) % input.size()];  // 下一顶点
        
        // 判断两个端点与边界的位置关系
        bool currentInside = IsInside(current, x1, y1, x2, y2, edge);
        bool nextInside = IsInside(next, x1, y1, x2, y2, edge);
        
        // 根据四种情况决定输出
        if (currentInside && nextInside) {
            // 情况1: 两点都在内侧 -> 输出终点
            output.push_back(next);
        } else if (currentInside && !nextInside) {
            // 情况2: 起点内侧,终点外侧 -> 输出交点
            D2D1_POINT_2F intersection = ComputeIntersection(
                current, next, x1, y1, x2, y2, edge);
            output.push_back(intersection);
        } else if (!currentInside && nextInside) {
            // 情况4: 起点外侧,终点内侧 -> 输出交点和终点
            D2D1_POINT_2F intersection = ComputeIntersection(
                current, next, x1, y1, x2, y2, edge);
            output.push_back(intersection);
            output.push_back(next);
        }
        // 情况3: 两点都在外侧 -> 不输出任何点
    }
    return output;
}
\end{lstlisting}

\end{tcolorbox}

\begin{tcolorbox}
\subsubsection{辅助函数实现}

\textbf{代码思路}：

点到线段距离函数用于图形的选中检测（HitTest），判断鼠标点击位置是否靠近图形边界。

\textbf{算法原理}：
\begin{enumerate}
    \item 将点$P$投影到线段$AB$所在直线上，得到投影点$P'$
    \item 投影参数$t = \frac{\vec{AP} \cdot \vec{AB}}{|\vec{AB}|^2}$
    \item 将$t$限制在$[0,1]$范围内，确保投影点在线段上
    \item 计算点$P$到投影点$P'$的欧氏距离
\end{enumerate}

\begin{lstlisting}[caption={点到线段距离计算}]
// ========== 点到线段距离 ==========
// 用于图形选中检测(HitTest)
// 输入: p - 测试点, a,b - 线段端点
// 输出: 点到线段的最短距离
inline float PointToLineDistance(D2D1_POINT_2F p, 
    D2D1_POINT_2F a, D2D1_POINT_2F b) {
    float dx = b.x - a.x;
    float dy = b.y - a.y;
    float len = sqrt(dx * dx + dy * dy);
    
    // 线段退化为点的情况
    if (len < 0.001f) 
        return sqrt((p.x-a.x)*(p.x-a.x) + (p.y-a.y)*(p.y-a.y));
    
    // 计算投影参数t,限制在[0,1]范围内
    float t = ((p.x - a.x) * dx + (p.y - a.y) * dy) / (len * len);
    t = max(0.0f, min(1.0f, t));
    
    // 计算投影点坐标
    float px = a.x + t * dx;
    float py = a.y + t * dy;
    
    // 返回点到投影点的距离
    return sqrt((p.x - px) * (p.x - px) + (p.y - py) * (p.y - py));
}
\end{lstlisting}

\subsubsection{多边形工具函数}

\textbf{代码思路}：

这组工具函数为多边形裁剪算法提供基础支持。

\textbf{1. 射线法判断点在多边形内部}：
\begin{itemize}
    \item 从测试点向右发射一条水平射线
    \item 统计射线与多边形边界的交点数
    \item 交点数为奇数则在内部，偶数则在外部
\end{itemize}

\textbf{2. 带符号面积计算}：
\begin{itemize}
    \item 使用向量叉积公式：$A = \frac{1}{2}\sum_{i=0}^{n-1}(x_i y_{i+1} - x_{i+1} y_i)$
    \item 面积为正表示逆时针，为负表示顺时针
    \item 可用于判断多边形的绑定方向
\end{itemize}

\begin{lstlisting}[caption={多边形工具函数}]
// 判断点是否在多边形内部（射线法）
static bool PointInPolygon(D2D1_POINT_2F point, 
    const std::vector<D2D1_POINT_2F>& polygon) {
    int count = 0;
    size_t n = polygon.size();
    
    for (size_t i = 0; i < n; i++) {
        D2D1_POINT_2F p1 = polygon[i];
        D2D1_POINT_2F p2 = polygon[(i + 1) % n];
        
        if ((p1.y <= point.y && p2.y > point.y) || 
            (p2.y <= point.y && p1.y > point.y)) {
            float x = p1.x + (point.y - p1.y) / 
                (p2.y - p1.y) * (p2.x - p1.x);
            if (point.x < x) count++;
        }
    }
    return (count % 2) == 1;
}

// 计算多边形面积（带符号，可判断方向）
static float SignedArea(const std::vector<D2D1_POINT_2F>& polygon) {
    float area = 0;
    size_t n = polygon.size();
    for (size_t i = 0; i < n; i++) {
        size_t j = (i + 1) % n;
        area += polygon[i].x * polygon[j].y;
        area -= polygon[j].x * polygon[i].y;
    }
    return area / 2.0f;
}

// 判断多边形是否为顺时针方向
static bool IsClockwise(const std::vector<D2D1_POINT_2F>& polygon) {
    return SignedArea(polygon) < 0;
}
\end{lstlisting}

\end{tcolorbox}

\begin{tcolorbox}
\subsubsection{Weiler-Atherton多边形裁剪实现}

\textbf{代码思路}：

Weiler-Atherton算法是最复杂但也最强大的多边形裁剪算法，能够处理凹多边形并输出多个独立区域。

\textbf{算法五个步骤}：

\textbf{步骤1 - 计算所有交点}：遍历主多边形和裁剪多边形的所有边对，计算交点坐标和参数位置。

\textbf{步骤2 - 判断进入/离开}：对每个交点，检查交点后一小段位置是否在裁剪区域内部：
\begin{itemize}
    \item 若在内部，则为"进入点"（从外到内）
    \item 若在外部，则为"离开点"（从内到外）
\end{itemize}

\textbf{步骤3 - 构建顶点链表}：将交点按参数$t$顺序插入到两个多边形的顶点链表中。

\textbf{步骤4 - 建立双向链接}：为每个交点建立指向另一个链表中对应交点的链接，用于遍历时切换。

\textbf{步骤5 - 遍历生成结果}：
\begin{itemize}
    \item 从一个未访问的进入点开始
    \item 沿主多边形方向行走，收集路径上的点
    \item 遇到离开点时，通过链接切换到裁剪多边形
    \item 沿裁剪多边形行走，直到遇到进入点
    \item 通过链接切换回主多边形，继续行走
    \item 回到起点时，完成一个裁剪区域
\end{itemize}

这种"交替行走"的策略使得算法能够正确处理凹多边形产生的多个独立区域。

\textbf{实现要点}：
\begin{itemize}
    \item 使用\texttt{Vertex}结构体存储顶点信息，包括坐标、参数、类型标记等
    \item \texttt{neighbor}字段实现两个链表之间的双向链接
    \item \texttt{visited}字段避免重复遍历同一个交点
    \item 判断进入/离开时，取交点后$\epsilon=0.001$处的点进行测试
    \item 结果可能包含多个独立的多边形区域
\end{itemize}

\begin{lstlisting}[caption={Weiler-Atherton多边形裁剪算法核心实现}]
// ========== Weiler-Atherton多边形裁剪算法 ==========
// 输入: subject - 主多边形(被裁剪), clipPoly - 裁剪多边形
// 输出: 裁剪后的多边形列表(可能有多个)
// 原理: 基于交点链表,从进入点遍历生成裁剪结果

// 顶点节点结构
struct Vertex {
    D2D1_POINT_2F pt;    // 顶点坐标
    float t;              // 交点在边上的参数位置
    bool isIntersection;  // 是否为交点
    bool isEntering;      // true=进入点, false=离开点
    bool visited;         // 是否已访问
    int neighbor;         // 对应另一链表中的索引
};

static std::vector<std::vector<D2D1_POINT_2F>> ClipPolygons(
    const std::vector<D2D1_POINT_2F>& subject,
    const std::vector<D2D1_POINT_2F>& clipPoly) {
    
    std::vector<std::vector<D2D1_POINT_2F>> result;
    
    // ===== 步骤1: 计算所有交点 =====
    std::vector<IntersectionInfo> intersections;
    for (size_t i = 0; i < subject.size(); i++) {
        D2D1_POINT_2F s1 = subject[i];
        D2D1_POINT_2F s2 = subject[(i + 1) % subject.size()];
        
        for (size_t j = 0; j < clipPoly.size(); j++) {
            D2D1_POINT_2F c1 = clipPoly[j];
            D2D1_POINT_2F c2 = clipPoly[(j + 1) % clipPoly.size()];
            
            float ts, tc; D2D1_POINT_2F pt;
            if (ComputeIntersection(s1, s2, c1, c2, ts, tc, pt)) {
                // 判断进入/离开: 检查交点后一小段是否在裁剪区域内
                D2D1_POINT_2F afterPt = {
                    s1.x + (ts + 0.001f) * (s2.x - s1.x),
                    s1.y + (ts + 0.001f) * (s2.y - s1.y)
                };
                bool entering = PointInsidePolygon(afterPt, clipPoly);
                intersections.push_back({pt, i, j, ts, tc, entering});
            }
        }
    }
    
    // ===== 步骤2: 构建主多边形顶点链表(插入交点) =====
    std::vector<Vertex> subjList;
    for (size_t i = 0; i < subject.size(); i++) {
        Vertex v; v.pt = subject[i]; v.isIntersection = false;
        subjList.push_back(v);
        // 收集该边上的交点,按参数t排序后插入
        // ... (按t_S升序插入交点)
    }
    
    // ===== 步骤3: 构建裁剪多边形顶点链表 =====
    std::vector<Vertex> clipList;
    // ... (同步骤2)
    
    // ===== 步骤4: 建立交点双向链接 =====
    for (size_t i = 0; i < subjList.size(); i++) {
        if (subjList[i].isIntersection) {
            for (size_t j = 0; j < clipList.size(); j++) {
                if (clipList[j].isIntersection) {
                    // 找到坐标相同的交点,建立链接
                    float dx = subjList[i].pt.x - clipList[j].pt.x;
                    float dy = subjList[i].pt.y - clipList[j].pt.y;
                    if (dx*dx + dy*dy < 1e-4f) {
                        subjList[i].neighbor = j;
                        clipList[j].neighbor = i;
                        break;
                    }
                }
            }
        }
    }
    
    // ===== 步骤5: 从进入点遍历生成裁剪结果 =====
    for (size_t start = 0; start < subjList.size(); start++) {
        // 只从未访问的进入点开始
        if (!subjList[start].isIntersection || 
            !subjList[start].isEntering || 
            subjList[start].visited) continue;
        
        std::vector<D2D1_POINT_2F> polygon;
        int current = start;
        bool onSubject = true;  // 当前在主多边形上
        
        do {
            if (onSubject) {
                Vertex& v = subjList[current];
                polygon.push_back(v.pt);
                v.visited = true;
                
                // 若是离开点,切换到裁剪多边形
                if (v.isIntersection && !v.isEntering && v.neighbor >= 0) {
                    current = (v.neighbor + 1) % clipList.size();
                    onSubject = false;
                } else {
                    current = (current + 1) % subjList.size();
                }
            } else {
                Vertex& v = clipList[current];
                // 若是进入点,切换回主多边形
                if (v.isIntersection && v.neighbor >= 0 && 
                    subjList[v.neighbor].isEntering) {
                    current = v.neighbor;
                    onSubject = true;
                } else {
                    polygon.push_back(v.pt);
                    current = (current + 1) % clipList.size();
                }
            }
        } while (!(onSubject && current == start));
        
        if (polygon.size() >= 3) result.push_back(polygon);
    }
    return result;
}
\end{lstlisting}

\end{tcolorbox}

\clearpage

\begin{tcolorbox}

\section{实验总结}

通过本次实验，深入理解了计算机图形学中图形变换和裁剪算法的原理和实现方法：

\begin{enumerate}
    \item \textbf{齐次坐标与变换矩阵}：使用齐次坐标可以将平移、旋转、缩放等变换统一用矩阵乘法表示，便于复合变换的计算。绑定任意点的旋转和缩放需要先平移到原点，变换后再平移回去。
    
    \item \textbf{Liang-Barsky算法}：该算法基于参数化直线方程，通过计算直线与窗口边界的交点参数来确定裁剪结果。相比Cohen-Sutherland算法，Liang-Barsky算法减少了交点计算次数，效率更高。
    
    \item \textbf{Sutherland-Hodgman算法}：该算法采用逐边裁剪的策略，依次对四条边界进行裁剪。算法简单易实现，但只能输出单个连续的多边形，对于产生多个独立区域的情况无法正确处理。
    
    \item \textbf{Weiler-Atherton算法}：该算法基于交点链表，可以处理凹多边形、自相交多边形和带孔洞的多边形，能够输出多个独立的多边形。算法较为复杂，但功能更加强大。
    
    \item \textbf{用户交互设计}：良好的用户交互设计可以大大提高软件的易用性。本实验实现了鼠标拖动平移、滚轮旋转/缩放、自定义中心点等交互方式，使用户操作更加直观便捷。
\end{enumerate}

本实验加深了对计算机图形学变换和裁剪算法的理解，提高了使用Direct2D进行图形编程的能力。
\end{tcolorbox}

\end{document}
